\section{Community Movement}
\sectionkicker{COMMUNITY}
\begin{wideflow}
{\Large\bfseries 公式発表の到着記録と、確認できなかった独立検証}\par
\smallskip
{\color{SurveyMuted}14件の公開投稿はすべて発行元告知であり、独立した技術検証の到着は確認できない。}\par
\end{wideflow}

公式発表は5領域に集中した。最前線の提供条件（GPT-6、Haiku 5.5）、数学成果の公開、安全の制度発足（Cyber Missionと基盤防衛・検査器）、公開ウェイトの予告（Mistral Large 4、Beam）、音声・画像の提供（Nano Banana 2.1、EmbeddingGemma 2）である\cite{gpt6ui}\cite{mathrelease}\cite{haiku55}\cite{cybermission}\cite{mistral4}\cite{beam}\cite{nanobanana}\cite{embedgemma2}。いずれも発行元による発表の到着であり、仕様や性能、提供条件の独立した裏付けではない。

14件の内訳は、通常期間内12件と締め切り後2件である\cite{gpt6ui}\cite{mathrelease}\cite{haiku55}\cite{cybermission}\cite{mistral4}\cite{beam}\cite{nanobanana}\cite{embedgemma2}\cite{clefomni}。件数の水増しはせず、リンク総数と期間内訳を分けて示す。行単位の台帳を正とし、要約文の数えに頼らない。締め切り後の2件（10月9日22時04分の動作報告、同日23時11分のClef-omni拡張告知）は通常の集計には含めない\cite{clefomni}。Clef-omniの公開時刻の根拠は発行元ブログの記録であり、投稿日時は公開日時の根拠にはならない。

\begin{communitynote}[Weekly Community Movement --- context only]
今号の公開投稿14件はすべて発行元公式であり、独立・一般利用者の技術投稿は0件である。公式発表の到着は確認できるが、独立した再現、導入確認、評価、批評の到着は確認できない。発表件数から技術コミュニティの支持や性能評価を推定できない。締め切り後の速報は通常の集計には含めない。
\end{communitynote}

独立した技術検証、採用報告、再現報告、批評は確認されていない。保存済みの公開投稿資料には利用者や第三者の反応（祝意、ベンチ開始の表明、利用可能性の報告など）が含まれるが、いずれも関心の表明であって、仕様や性能、提供条件の検証ではない。重みの入手確認や局所実行の報告、評価の再実行など、高信号の独立確認は見当たらない。UltrafastとDecisions、Ironcladに対応する発行元投稿は期間内に見当たらず、手引きや研究ページの確認が根拠になる\cite{ultrafast}\cite{decisions}\cite{ironclad}。Gemini仕事エージェントに対応する一次投稿は見当たらない\cite{geminiagent}。VercelやPixel Canaryは期間前の文脈のみであり、日本語特化の公式投稿は第三者報道のみで公式到着はない。音声・映像生成の新規公開、serving系の主版更新も期間内の一次到着はない。

発表件数から技術コミュニティの支持や性能評価を推定できない。公式投稿が多いことは、発行元が発表したという事実を示すだけで、独立した利用者が試し、測り、採用したことを示さない。Mistralの49B表記と52B文書の差は、告知の存在と文書記載の差であって、重み実体の確認ではない\cite{mistral4}。Nano Bananaに関する利用可能性の報告は、新規公開の告知ではない\cite{nanobanana}。Beamへの祝意表明は反応であって、仕様の裏付けではない\cite{beam}。いずれも一次発表として本文で明示的に引用した公式投稿の発表内容に限って一次資料として扱い、台帳に載っているという事実だけでは技術的な根拠に格上げしない。

今後確認すべきは、一次的な実装・評価記録である。重み実体の取得確認と許諾条項、局所実行や推論基盤での動作報告、評価の独立した再実行、料金の請求実態、検出器や検査器の第三者評価、技術報告PDFの取得後の精査が次の観測点になる。根拠が足りない動向は無理に記事化せず、限定的に示すにとどめる。全行の内訳は本号の公開台帳で直接確かめられる。
